\subsection{The exact-environment baseline fails to solve the gate}
\label{sec:results-ppo}
\label{sec:results}

None of the 15 policies trained with exact transitions and true task signals reaches $50\%$ survival.
Their mean survival is $0.47\%$, $0.39\%$, and $1.41\%$ for model seeds 0, 1, and 2
(Figure~\ref{fig:matched-control}; Table~\ref{tab:matched-gate} in the appendix). The same controller is therefore unsuccessful even when transition,
reward, and termination prediction errors are removed from its training environment. The poor policy performance in imagination therefore cannot be attributed to the learned transition
model alone.

\paragraph{A fixed training protocol.}
For every frozen model, we train five recurrent MLP-PPO policies \citep{schulman2017ppo} in each of
three conditions: learned transitions with fitted reward/termination heads; exact simulator
transitions re-encoded by the same encoder with the same fitted heads; and exact re-encoded
transitions with true rewards and terminations. Actor inputs contain no simulator state. The 45 runs
share their start-state mixture, context burn-in, rollout length, optimizer, updates, and evaluation
keys within model/policy-seed pairs. Each completes 150 updates of $512\times16$ transitions and is
evaluated at its final checkpoint on 256 real episodes. Walls passed counts exits before or on
the terminal transition, including an exit on the same step as death. Survival requires remaining
alive through step 128; these are the two primary outcomes;
survival above $50\%$ is a secondary operational threshold.

Exact-simulator trajectories become absorbing after true termination. With fitted heads, learned
continuation may still weight those terminal states. We report this support change alongside the
transition intervention. Policy-seed intervals are conditional on a fixed model and its evaluation
set; the three model effects remain separate. Appendix~\ref{app:matched} records the pairing checks,
uncertainty estimands, and all contrasts.

Exact transitions and true signals nevertheless improve walls passed from $0.285$ to $0.452$,
$0.217$ to $0.341$, and $0.320$ to $0.543$ in the three models. The paired gains are $0.166$,
$0.123$, and $0.223$ walls per episode. Survival changes are $-0.47$, $-0.39$, and $+1.02$
percentage points. Thus the substitution changes partial progress, while leaving task-level control
poor in every fitted model. All 45 policies complete the prescribed budget; no intermediate policy
or favorable model seed is selected.

\begin{figure}[t]
\centering
\paperfig{matched_control_v2.pdf}{1.0}
\caption{\textbf{The exact-environment control is itself unsuccessful.} All 45 final policies,
with five paired policy seeds per model and training condition. Lines connect matched policy seeds;
model fits remain separate. The survival panels use a $0$--$4\%$ detail scale, far below the $50\%$
operational threshold. Exact+heads$^*$ uses absorbing true terminal states while retaining fitted
reward and continuation, so this transition substitution also changes terminal support.}
\label{fig:matched-control}
\end{figure}

\subsection{Transition and task-signal substitutions have different effects}
\label{sec:results-mean}

Keeping the fitted heads while substituting exact transitions lowers mean walls passed to $0.230$,
$0.068$, and $0.096$. Replacing those heads with true signals then raises the means to $0.452$,
$0.341$, and $0.543$. These contrasts expose sensitivity to how transitions and task signals are
combined. They do not establish that less accurate dynamics are intrinsically better: the fitted-head
condition applies its learned continuation after the simulator has entered an absorbing terminal state.
Its logged post-terminal weight fraction averages $12.3\%$, $11.2\%$, and $9.4\%$ across updates and
policy seeds in the three models.

Real-transition head diagnostics already show a weak wall-collision signal. For model 0, fitted
reward is about $-0.05$ on wall deaths and $-0.89$ on border deaths, against the true $-1$;
termination AUC is $0.65$ versus $0.96$. Yet correcting both dynamics and signals still leaves the
common MLP-PPO recipe unsuccessful. The exact condition leaves the actor's access to encoded state and its training procedure as
unresolved limitations.

\subsection{A native spatial policy also leaves the task unsolved}
\label{sec:native-grid}
A separate sensitivity replaces the MLP with a CNN on the native $4\times4\times32$ latent grid,
without a decoder, extra latent coordinates, or privileged inputs. It repeats all three training
conditions for three policy seeds per model, adding 27 complete runs. None exceeds $50\%$ survival.
With exact transitions and true signals, mean survival is $1.17\%$, $0.52\%$, and $0.65\%$;
wall counts are $0.616$, $0.367$, and $0.452$ (Appendix~\ref{app:native-grid}). Thus changing this
policy inductive bias does not qualify the short-rollout recipe as a successful control test.

The transition substitution has a different sign under this torso: exact transitions with fitted
heads increase mean wall counts in all three models, whereas they decrease them in the primary MLP
study. We retain the two campaigns separately. Across both, none of 72 policies crosses the threshold,
including all 24 trained with exact transitions and true signals. Earlier real-environment CNNs on
supervised decoded encodings do cross it in two of nine runs, using a different training protocol;
those results establish a separate useful read-out, not a successful common baseline
(Appendix~\ref{app:historical-gate}).


\subsection{Correct termination removes the extension but does not resolve control}
\label{sec:termination}
The observed post-terminal weights trigger a separate sensitivity with exact transitions, true
termination, and the original learned reward head. All 24 additional policies complete the fixed
budget and have exactly zero continuation-weighted post-terminal mass. None exceeds $50\%$ survival.
Changing termination alone has mixed effects on walls passed and lowers the mean in all three
native-grid cells (Appendix~\ref{app:termination}).

With exact transitions and true termination held fixed, replacing the fitted reward with the true
reward increases mean wall counts in all six model/torso cells: $0.287$, $0.124$, and $0.465$ walls
for the MLP, and $0.514$, $0.367$, and $0.271$ for the native-grid CNN. Survival changes remain small
and mixed. This identifies a contribution from the fitted reward interface within these training
conditions, while the unsuccessful exact-condition controls leave encoder accessibility, policy
interface, and training procedure unresolved. The separate horizon-16 campaigns total 96 policies
with no threshold crossings, all conditioned on only three fitted world models.
